\section{Conclusion}
\label{sec:conclusion}

Timing-integrity monitors are plentiful and each publishes its own numbers on its
own terms; what the field lacks is a common ruler. This paper defines that ruler:
an open, versioned, cross-protocol conformance benchmark that scores any monitor's
stated integrity risk against its measured undetected-error exceedance, under one
fault menu, with the Stanford integrity diagram specialized to time as its scoring
method (\Cited{}), a parametric ten-class fault menu as its stimulus (\Modelled{}),
and a four-value verdict algebra as its core. The verdict algebra is the load-bearing
contribution: because its enumeration contains no detection token, an undetectable
fault---a symmetric single-path delay or a replay within freshness---can only be
scored as \Absorbed{} or \Unabsorbed{} by the protection level, never as detected.
Mizrahi's impossibility is thereby made a machine-checked, type-level invariant
(\Proven{}), and a negative result is turned into a positive scoring rule: a
monitor is rewarded for carrying enough protection level to absorb what it cannot
see, never for the impossible act of detecting it.

The benchmark's defining property is that it makes no accuracy claim of its own. It
does not assert that any clock is accurate, any protection level tight, or any
monitor safe; it applies a fixed rule to a fixed menu and tallies the outcome.
Because it asserts nothing about accuracy, it cannot be wrong about accuracy---it is
honesty-immune---and so it carries \emph{zero} \Validated{} tags by construction,
not by omission. The single cross-check we run, an independent re-count of the
region tallies, verifies the scorer's arithmetic and is labelled \Modelled{}~/~
\code{InternalConsistency}: it checks the counting, not any accuracy, and is
explicitly neither validation nor an external dataset. This is what makes the object
worth citing---a reference that cannot overclaim, reproducible by one command.

We demonstrated the benchmark by scoring the companion composed protection level
and a deliberately broken monitor. The reference monitor covered its holdover
design regime, absorbed both sub-floor undetectable offsets, and---crucially---was
\emph{honestly exposed} at its boundary, drawing hazardously misleading information
on the adversarial faults outside its model ($n_{\HMI{}}^{\text{total}}=64$); the
broken monitor was strictly worse on the same rule
($n_{\HMI{}}^{\text{total}}=121$, both undetectables unabsorbed). A reference
yardstick maps a method's coverage rather than flattering it. This benchmark is the
yardstick half of a citable pair whose method half is the companion composed timing
protection level \citep{baweja2026cti}; the adversarial and common-mode papers of
the arc extend the monitors it scores. What this paper contributes on its own is
small and, we hope, durable: a common ruler the field did not have, with its
central impossibility compiled into the shape of the result, and every number
reproducible from the open engine.
